\section*{第 \ref{ch:g9:thales} 章——泰勒斯定理}

\begin{solution}{exo:g9:thales:1}
由泰勒斯定理，$\dfrac{AM}{AB} = \dfrac{AN}{AC} = \dfrac{MN}{BC}$，即\
$\dfrac26 = \dfrac13$。于是 $\dfrac{3}{AC} = \dfrac13$ 得 $AC = 9$，
由 $\dfrac{MN}{9} = \dfrac13$ 得 $MN = 3$。
\end{solution}

\begin{solution}{exo:g9:thales:2}
先求 $AB = AM + MB = 5 + 3 = 8$。然后
由 $\dfrac{AM}{AB} = \dfrac{AN}{AC}$ 得 $\dfrac58 = \dfrac{4}{AC}$，于是
$5 \times AC = 32$，故 $AC = 6.4$。
\end{solution}

\begin{solution}{exo:g9:thales:3}
蝴蝶形的情形与嵌套形完全一样：
$\dfrac{AM}{AB} = \dfrac{AN}{AC} = \dfrac{MN}{BC}$，即\
$\dfrac{3}{7.5} = \dfrac25$。于是 $\dfrac{2}{AC} = \dfrac25$ 得
$AC = 5$，又由 $\dfrac{2.6}{BC} = \dfrac25$ 得
$BC = \dfrac{2.6 \times 5}{2} = 6.5$。
\end{solution}

\begin{solution}{exo:g9:thales:4}
$\dfrac{AM}{AB} = \dfrac68 = \dfrac34$，而
$\dfrac{AN}{AC} = \dfrac{9}{12} = \dfrac34$：两比值相等，且对应点的
顺序一致，由泰勒斯定理的逆定理得 $(MN) \parallel (BC)$。
\end{solution}

\begin{solution}{exo:g9:thales:5}
$\dfrac{AM}{AB} = \dfrac46 = \dfrac23$，而
$\dfrac{AN}{AC} = \dfrac58$。验算：$\dfrac23 = \dfrac{16}{24}$
且 $\dfrac58 = \dfrac{15}{24}$，两比值不相等。若两直线
\omterm{def:g6:lines:perp}{平行}，则由泰勒斯定理两比值必相等，故 $(MN)$
与 $(BC)$ 并\emph{不}\omterm{def:g6:lines:perp}{平行}。
\end{solution}

\begin{solution}{exo:g9:thales:6}
设 $S$ 为影子的顶端，$P$ 为人的头顶，$L$
为灯的顶端。人（高 $1.8$ m，距 $3$ m，以 $S$ 为基准）与
灯（高 $h$，距 $3 + 2 = 5$ m，以 $S$ 为基准）是两条
\omterm{def:g6:lines:perp}{平行}竖直\omterm{def:g6:lines:objects}{线段}：被光线 $(SL)$ 所截。
由点 $S$ 出发，应用泰勒斯定理得
\[
\frac{1.8}{h} = \frac{3}{5},
\qquad\text{故}\quad
h = \frac{1.8 \times 5}{3} = 3 \text{ m}.
\]
\end{solution}

\begin{solution}{exo:g9:thales:7}
\emph{1.} 实际距离：$6.8 \times 25\,000 = 170\,000$ cm $=
1.7$ km。

\emph{2.} \omterm{def:g6:measure:area}{面积}按 $k^2$ 缩放：实际\omterm{def:g6:measure:area}{面积} $= 8 \times 25\,000^2$
cm$^2$ $= 8 \times 6.25 \times 10^8$ cm$^2$ $= 5 \times 10^9$ cm$^2$，
由于 $1$ km$^2 = 10^{10}$ cm$^2$，所以实际\omterm{def:g6:measure:area}{面积}为 $0.5$ km$^2$。
\end{solution}

\begin{solution}{exo:g9:thales:8}
\emph{1.} 比值为 $\dfrac{AM}{AB} = \dfrac{8}{12} = \dfrac23$。
由泰勒斯定理得 $AN = \frac23 \times 15 = 10$，且 $MN = \frac23 \times 18 = 12$。

\emph{2.} 相似比为 $k = \frac23$。\omterm{def:g6:measure:perimeter}{周长}方面，$ABC$ 的\omterm{def:g6:measure:perimeter}{周长}为
$12 + 15 + 18 = 45$，而 $AMN$ 的\omterm{def:g6:measure:perimeter}{周长}为 $8 + 10 + 12 = 30 = \frac23 \times
45$：\omterm{def:g6:measure:perimeter}{周长}同样按 $k$ 缩放，与所有长度一样。
\end{solution}

\begin{solution}{exo:g9:thales:9}
设 $O$ 为直线 $(AC)$ 与 $(BD)$ 的交点，且
$(AB) \parallel (CD)$：由三角形 $OAB$
与 $OCD$ 构成蝴蝶形。由泰勒斯定理得
\[
\frac{OA}{OC} = \frac{OB}{OD} = \frac{AB}{CD} = \frac46 = \frac23 .
\]
于是 $\dfrac{OA}{OC} = \dfrac23$，而等式
$\dfrac{OA}{OC} = \dfrac{OB}{OD}$ 恰说明 $O$ 把两条对角线
按相同的比分割。
\end{solution}

\begin{solution}{pb:g9:thales:1}
\textbf{1.} 作图得\omterm{def:g6:lines:objects}{线段} $[AB]$ 上的两点；
记为 $M_1$（由 $P_1$ 得）与 $M_2$（由 $P_2$ 得）。

\textbf{2.} 在三角形 $A P_3 B$ 中，\omterm{def:g6:lines:perp}{平行}于 $(P_3 B)$ 且过 $P_1$、$P_2$ 的直线
截得两边成比例
（\cref{thm:g9:thales:direct}）：
\[
\frac{A M_1}{AB} = \frac{A P_1}{A P_3} = \frac13,
\qquad
\frac{A M_2}{AB} = \frac{A P_2}{A P_3} = \frac23 .
\]
圆规所取的 $A P_1$、$P_1 P_2$、$P_2 P_3$ 相等，故
各比值恰为三分之一——无论射线的角度
与圆规张口如何选取。

\textbf{3.} 五等分：在射线上截取五段相等长度，把第五个分点
与 $B$ 相连，并作四条\omterm{def:g6:lines:perp}{平行}线。求\omterm{def:g6:lines:objects}{线段}的 $\frac35$：
同一图形，取过 $P_3$ 的\omterm{def:g6:lines:perp}{平行}线截得的点，
它即 $\frac{3}{5}$ 分点，位于 $[AB]$ 上距 $A$ 处。

\textbf{4.} 由 $\frac{AM}{MB} = \frac23$ 知
$AM = \frac25 AB$：在射线上取五段相等长度，把 $P_5$ 与
$B$ 相连，过 $P_2$ 的\omterm{def:g6:lines:perp}{平行}线即定出 $M$。

\textbf{5.} 度量 $\sqrt2 = 1.41421\ldots$ cm 时，
只能取有限位\omterm{def:g6:decimals:places}{小数}：任何量得的“三分之一”都
只是近似值。泰勒斯作图从不读取数值：它
给出的正是 $\frac{\sqrt2}{3}$ 处的精确点——这种\omterm{def:g6:lines:objects}{线段}的等分，
连直尺也无法表达
（\cref{pb:g9:sqrt:1}）。

\textbf{6.} 物体顶端、小孔与像的底端
三点共线，物体底端与像顶端同样共线，
两条直线相交于小孔，物体与
像所在墙面\omterm{def:g6:lines:perp}{平行}。小孔成像中的泰勒斯比例：
\[
\frac{h}{H} = \frac{d}{D},
\qquad\text{故}\qquad
h = H \times \frac dD .
\]

\textbf{7.} $h = 12 \times \frac{0.20}{30} = 0.08$ m $= 8$ cm
——倒立（光线在小孔处相交）。

\textbf{8.} $H = h \times \frac Dd = 9 \times 10^{-3} \times
\frac{1.5 \times 10^{11}}{1} = 1.35 \times 10^{9}$ m。真实
值 $1.39 \times 10^9$ m：一只鞋盒测太阳，
误差仅为百分之几。

\textbf{9.} 月球：$\frac{3.5 \times 10^6}{3.8 \times 10^8}
\approx 0.0092$。太阳：$\frac{1.39 \times 10^9}{1.5 \times
10^{11}} \approx 0.0093$。两个比值——即
天空中的视大小——几乎完全相等：月球能够遮住太阳，
\emph{恰好}，边缘对边缘。这一宇宙中的巧合便是
日全食。

\textbf{10.} 每条光线都必须穿过同一小孔，因此
物体顶部发出的光线\emph{向下}延伸，底部发出的
光线\emph{向上}延伸：像是倒立的。扩大小孔
会让整束略有偏移的像的副本通过，
它们相互重叠：更亮，但变模糊。

\textbf{11.} $\frac{PQ}{MN} = \frac{OP}{OM} = \frac{7.5}{3}
= 2.5$，故 $PQ = 10$；且 $\frac{ON}{OQ} = \frac{OM}{OP} =
\frac{1}{2.5}$，故 $ON = \frac{6}{2.5} = 2.4$。

\textbf{12.} 若\omterm{def:g6:lines:perp}{平行}线经过某一边的中点，则泰勒斯比例为 $\frac12$，
因此它在\emph{另一边}的中点处相交，
且\omterm{def:g6:lines:perp}{平行}\omterm{def:g6:lines:objects}{线段}长为底边的\omterm{def:g3:division:half}{一半}：
即\cref{thm:g8:midpoints:direct}。泰勒斯定理是
脱离了比值 $\frac12$ 的中点定理。

\textbf{13.} 在三角形 $ACD$ 中，直线 $(EO)$ \omterm{def:g6:lines:perp}{平行}于
底边 $(CD)$，且 $\frac{AO}{AC} = \frac{OA}{OA + OC} =
\frac{4}{4 + 6} = \frac25$。泰勒斯比例：
$\frac{EO}{DC} = \frac{AO}{AC} = \frac25$，故
$EO = \frac25 \times 6 = 2.4$.

\textbf{14.} 在三角形 $BDC$ 中，同理
$\frac{BO}{BD} = \frac25$，且 $OF = \frac25 \times 6 = 2.4$。
由此
\[
EF = EO + OF = 4.8
= \frac{2 \times 4 \times 6}{4 + 6} :
\]
两底的\omterm{pb:g8:speed:1}{调和平均数}——正是往返行程的平均
（见\cref{pb:g8:speed:1}），在梯形中作出。

\textbf{15.} 算术平均数：$5$；\omterm{pb:g8:cosine:1}{几何平均数}：
$\sqrt{24} \approx 4.90$；\omterm{pb:g8:speed:1}{调和平均数}：$4.8$。顺序：
$4.8 < 4.90\ldots < 5$，即\omterm{pb:g8:speed:1}{调和平均数} $<$ \omterm{pb:g8:cosine:1}{几何平均数} $<$
算术平均数——仅当两底相等时取等号。在
图中：中位线长为 $5$，过
对角线交点的\omterm{def:g6:lines:perp}{平行}线长为 $4.8$，相似截线
为 $\sqrt{24}$，三者按此顺序排列在两底之间。
这些不等式已在\cref{pb:g8:cosine:1}中证明（\omterm{pb:g8:cosine:1}{几何平均数}
$\leq$ 算术平均数，借圆与恒等式证明），以及
\cref{pb:g8:speed:1}（\omterm{pb:g8:speed:1}{调和平均数} $\leq$ 算术平均数）；\omterm{pb:g8:speed:1}{调和平均数}
$\leq$ \omterm{pb:g8:cosine:1}{几何平均数}由二者结合得到（$H \times A = G^2$，
同样在该处证明）。
\end{solution}
